% ===========================================================================
% Section 3 -- Initial implementation
% ===========================================================================

\subsection{Tools and their choice}

\begin{itemize}
  \item \textbf{Language: Python 3.12.} The bottleneck is the MILP solver,
    not model construction: building the largest model we attempt (4 million
    non-zeros) takes a few seconds. Python lets the generator, the models,
    the validator, the experiment runner and the plots share one code base.
  \item \textbf{Solver: HiGHS 1.16}~\cite{huangfu2018} through its Python
    interface \texttt{highspy}. HiGHS is a state-of-the-art open-source
    MILP solver (branch-and-cut with presolve, cutting planes and primal
    heuristics). It needs no licence, so all results can be reproduced
    exactly by anyone, including the lecturer.
  \item \textbf{Second solver: Gurobi 14} through \texttt{gurobipy}. The same
    code can pass the models to Gurobi, which we use to cross-check the
    results on small instances. A full academic licence was not available
    when the baseline was run: the free licence is limited to 2000 variables
    and constraints, which covers only the verification instances of
    Section~\ref{sec:verification}. A comparison of solvers on the full
    suite is a possible improvement for Section~\ref{sec:improvements}.
\end{itemize}

To guarantee that both solvers receive \emph{exactly} the same model, the
models are not written in a solver-specific modelling language. Each builder
fills a solver-independent container (cost vector, bounds, integrality, and
the constraint matrix stored row-wise in sparse form) that is then passed to
either back-end (\texttt{src/milp.py}).

\subsection{Implementation of the models}

The three compact formulations of Section~\ref{sec:compact} are implemented
in \texttt{src/models/baseline.py}, exactly as written there, with no
preprocessing, symmetry breaking or valid inequalities (these are
improvements, Section~\ref{sec:improvements}):
\begin{itemize}[nosep]
  \item \textbf{(AF)}: one binary per employee and augmented arc;
    constraints~\eqref{af:cover}--\eqref{af:spread}.
  \item \textbf{(CF)}: binaries $y_k$, $x_{ik}$;
    constraints~\eqref{cf:cover}--\eqref{cf:work}, with the conflict set
    $C_\tau$ (time conflicts and pairwise spread conflicts) computed for each
    type.
  \item \textbf{(CF-M)}: as (CF) with only the time conflicts, plus the
    big-$M$ spread constraints~\eqref{cfm} with $M=H$.
\end{itemize}

\paragraph{Employee pool.} The pool sizes are computed as in
Section~\ref{sec:compact} from the cost $\bar z$ of a greedy solution:
$U_{\mathrm F}=\min(n,\lfloor\bar z\rfloor)$ and
$U_{\mathrm H}=\min(n_{\mathrm H},\lfloor\bar z/0.6\rfloor)$, where
$n_{\mathrm H}$ is the number of duties with $d_i\le W/2$. The greedy
(\texttt{src/heuristics/greedy.py}) scans duties by start time and appends
each to the compatible open shift that ended last, opening a new shift when
none fits. It is run once with full-time and once with half-time limits, and
the cheaper result is kept. For (CF) and (CF-M) the greedy only builds
pairwise-compatible shifts. This matters on non-transitive graphs: a pool
sized from a solution that (CF) cannot represent may be too small for (CF),
which is a restriction there (Remark after Proposition~\ref{prop:cf}). We
found this case while testing (Section~\ref{sec:verification}). The greedy
solution is \emph{not} given to the solver as a starting point: the baseline
uses the solver alone.

\paragraph{Size limit.} Models with more than 4 million non-zeros are not
built and are recorded as \emph{too large}. The memory of our machine
(16\,GB, shared by parallel runs) would not allow solving them anyway, and
the limit itself shows how far each formulation scales.

\subsection{Verification}
\label{sec:verification}

\paragraph{Independent validator.} Every solution returned by a solver is
decoded into a list of shifts and checked by \texttt{src/validator.py},
which uses only the instance data and none of the models. It checks that each
duty is covered exactly once, that consecutive duties of a shift are joined
by an arc of $G$, the working-time and total-time limits of the shift's type,
and that the recomputed cost equals the objective reported by the solver.

\paragraph{Hand-crafted instances.} Table~\ref{tab:tests} lists small
instances whose optimal cost we derived by hand; each targets one rule of the
problem. All three models, with both HiGHS and Gurobi, return the expected
optimum, and the validator accepts all their solutions (54 checks). Five
further tests confirm that the validator \emph{rejects} broken solutions: an
HT shift over its working time, a missing duty, a duplicated duty, a wrong
reported cost, and a shift that is not a path in $G$
(\texttt{tests/test\_small.py}).

\begin{table}[ht]
\centering
\small
\caption{Verification instances ($W=480$, $T=720$ unless stated; duty
intervals in minutes; all time-feasible pairs are arcs unless stated).}
\label{tab:tests}
\begin{tabular}{@{}lp{5.0cm}p{5.0cm}r@{}}
\toprule
Name & Duties & What it tests & Optimum \\
\midrule
all\_ht & $[0,100],[110,200]$ & both fit one HT shift & 0.6 \\
overlap & $[0,100],[50,150]$ & overlapping duties & 1.2 \\
work\_binds & $[0,200],[210,400]$ & work $390>W/2$: one FT beats two HT & 1.0 \\
ht\_spread & $[0,60],[600,660]$ & spread $660>T/2$: one FT & 1.0 \\
ft\_spread & $[0,60],[700,760]$ & spread $760>T$: two HT & 1.2 \\
mixed & $[0,200],[210,400],[50,150]$, arc $0\to1$ only & FT $\{0,1\}$ + HT $\{2\}$ & 1.6 \\
long\_duty & $[0,300]$ & $d>W/2$: FT only & 1.0 \\
3-partition & sizes $26,34,40,30,33,37$ back to back, $W=100$ & reduction of
  Proposition~\ref{prop:np} ($m=2$) & 2.0 \\
non\_transitive & $[0,50],[60,110],[120,170]$, arcs $0\to1\to2$ &
  (AF) exact: $0.6$; (CF), (CF-M) restriction: $1.2$ & 0.6 \\
\bottomrule
\end{tabular}
\end{table}

\subsection{Experimental setup}
\label{sec:setup}

All runs use HiGHS with \emph{default} parameters (relative gap tolerance
$10^{-4}$, random seed 0), one thread per run and a time limit of
\textbf{300\,s} for the branch-and-cut. Five runs are executed in parallel on
an Intel Core i5-12500H (12 cores) with 16\,GB of RAM under Windows~11; all
models run under the same conditions. For each run we record
(\texttt{src/run\_experiments.py}, one CSV row per run):
\begin{itemize}[nosep]
  \item model size (variables, constraints, non-zeros) and build time;
  \item the \emph{LP bound}: optimum of the LP relaxation, solved separately
    with a 60\,s limit, which measures the strength of the formulation;
  \item the \emph{root bound}: dual bound at the end of the root node, after
    presolve and HiGHS's cuts, read through a solver callback;
  \item final status, best objective, best bound, gap, number of nodes and
    running time;
  \item the number of FT and HT employees and the validator verdict.
\end{itemize}
Gaps are measured against the lower bound $\max(0.6\rho(G),\sum_i d_i/W)$ of
Section~\ref{sec:inst-stats} and against the greedy upper bound $\bar z$.
